% Chapter 6

\chapter{结论与建议}

本文通过构建多个Agent投资者，以88个因子数据作为输入，MLP模型作为Agent的决策模型，
二次效用函数作为奖励函数，PPO算法作为强化学习训练方式，探究了Agent投资者的预期共识定价效应。
实验的结果表明，投资者的预期共识度越高，股票的未来收益越高，这一结论在一系列稳健性检验后依然成立。
进一步的，本文从融资和融券约束的角度，探讨了预期共识定价效应的形成机制。
实验的结果表明，融资和融券约束是预期共识定价效应的重要形成机制，
在放松融资和融券约束后，预期共识定价效应显著减弱。
最后，本文从企业市值、股权异质性和市场周期的角度，分析了预期共识定价效应的异质性。
实验的结果表明，小市值企业、非国企企业和牛市周期中的企业，其预期共识定价效应显著更高。
上述研究揭示了在A股市场中，投资者的预期形成机制和市场定价机制，对于理解金融市场的运行规律和稳定性有着重要的意义。
本文对政策和研究的启示如下：

第一，对投资者而言，达成预期共识的证券意味着更高的未来收益，
因此，投资者应该关注其他投资者的预期，通过分散化资产，避免买入分歧较大的证券，
以获得更高的未来收益。

第二，对于政策制定者而言，需要关注市场中投资者的预期，做到良性引导，
避免形成预期共识的证券被过度炒作，导致市场泡沫化。
同时，通过完善融资融券制度，
允许投资者进行融资和融券，可以有效降低投资者的预期分歧，
引导资产正确定价。

最后，对于研究者和金融监管机构而言，应当重视人工智能在金融领域的应用，
通过构建投资者的方法，来系统地研究微观投资者的行为如何传导至宏观价格。
在市场监管、信息披露等领域，也应当重视人工智能技术，提高监管效率和信息披露质量。
